\section{Discussion}
\textbf{Developing survival analysis models capable of handling the heterogeneous complexities of CKD progression prediction.} Learning with missing modalities represents an active area of research \citep{wu2024deepmultimodallearningmissing_survey}, particularly when features exhibit varying dimensionalities and structural characteristics. Despite these challenges, incorporating diverse modalities has demonstrated clear benefits, with studies showing improved performance in clinical prediction tasks including mortality and readmission forecasting \citep{wu2024multimodal_missing_modalities}. Adapting these architectures—originally designed for handling missing modalities—to time-series survival analysis represents a critical advancement toward building more accurate and clinically realistic survival models. 

\textbf{Inclusion of competing events.} There are many competing conflicts when it comes to predicting ESRD such as mortality \citep{hsu2017statistical_competing_events_CKD}. While  \citep{li2023machine_classification} constructs a classification benchmark to predict the mortality of CKD patients, we note that it is crucial to consider both events in modeling as ESRD influences many other CKD-related complications, including mortality. Integrating these predicted complications would produce more realistic progression models where outcomes are highly multifaceted. 

\textbf{Multimodal Progression Models.} MIMIC-IV \citep{johnson2023mimic_ehr} contains substantially more features than our current benchmark incorporates, as shown by \citep{dana2024integratedmachinelearningsurvival_feature_analysis}. Future work should systematically explore these features to understand their predictive value. Underexplored variables like social determinants of health (SDOH) \citep{chen2020social_SDOH} are particularly promising, as they may reveal critical differences in disease progression between patients who develop ESRD and those with stable kidney function.

Building robust healthcare predictive systems requires multimodal models that integrate diverse feature sets—including images \citep{johnson2019mimiccxrjpglargepubliclyavailable}, text \citep{johnson2023mimic_ehr}, physiological signals \citep{jing2018rapid_physiological_signals}, and demographic data \citep{pollard2018eicu_demographic}—regardless of structural differences. Current CKD progression models typically ignore these features or rely on unrealistic imputation assumptions \citep{swain2023robust_classifier_imputation_CKD}. Comprehensive CKD progression modeling ultimately requires features spanning multiple modalities.

 